% title: Eventually Constant GCD of Shifted Power Sums
% short-title: Eventually Constant GCD
% abstract: We formalize the IMO 2024 Shortlist problem asking for all positive integers \(a,b\) such that \(\text{gcd}(a^n+b,\, b^n+a)\) is eventually constant. The proof first shows that any eventual constant \(g\) must divide \(2\text{gcd}(a,b)\), then normalizes by \(d=\text{gcd}(a,b)\) and uses Euler's theorem modulo \(d^2xy+1\) to force \(d=x=y=1\).
% subjclass: 03B35
% keywords: olympiad number theory, gcd, Euler theorem, Lean 4
% author: Mario Hernández M.

\subsection{Eventually Constant GCD of Shifted Power Sums}

\begin{problemstatement}
Determine all positive integers \(a\) and \(b\) such that there exists a positive integer \(g\) satisfying
\[
  \text{gcd}(a^n+b,\; b^n+a)=g
\]
for all sufficiently large positive integers \(n\).
\end{problemstatement}

We formalize the predicate
\lean{Biblioteca.Demonstrations.EventuallyConstantGCD} and prove the
classification theorem
\lean{Biblioteca.Demonstrations.eventually_constant_gcd_classification}.

\begin{theorem}
The only pair \((a,b)\) of positive integers for which there exists a positive
integer \(g\) such that
\[
  \text{gcd}(a^n+b,\; b^n+a)=g
\]
for all sufficiently large \(n\) is \((a,b)=(1,1)\).
\end{theorem}

\begin{proof}
The pair \((1,1)\) clearly works, since
\[
  \text{gcd}(1^n+1,\;1^n+1)=2
\]
for every \(n\).

Conversely, assume that for some positive integers \(a,b,g\) and all
sufficiently large \(n\) we have
\[
  \text{gcd}(a^n+b,\;b^n+a)=g.
\]
Fix such a large \(N\). Since \(g\mid a^N+b\) and \(g\mid a^{N+1}+b\), we also
have
\[
  g \mid a(a^N+b)-(a^{N+1}+b)=b(a-1).
\]
Similarly,
\[
  g \mid a(b-1).
\]
From these two divisibilities one obtains \(ab\equiv a\equiv b\pmod g\), hence
\[
  a \equiv b \pmod g
\]
and then
\[
  a^2 \equiv a \pmod g.
\]
Therefore every positive power of \(a\) is congruent to \(a\) modulo \(g\). In
particular, using again \(g\mid a^N+b\), we get
\[
  a+b \equiv 0 \pmod g.
\]
Combining this with \(a\equiv b \pmod g\) yields
\[
  g \mid 2a
  \qquad\text{and}\qquad
  g \mid 2b.
\]
So if \(d=\text{gcd}(a,b)\), then \(g\mid 2d\). On the other hand, \(d\) divides both
\(a^N+b\) and \(b^N+a\), hence \(d\mid g\). Thus \(g=d e\) for some divisor
\(e\) of \(2\), so \(e\le 2\).

Now write
\[
  a=dx,\qquad b=dy,
\]
where \(x\) and \(y\) are coprime. Then for every sufficiently large \(n\),
\[
  \text{gcd}(a^n+b,\;b^n+a)
  = d\,\text{gcd}(d^{n-1}x^n+y,\; d^{n-1}y^n+x)
  = d e,
\]
and therefore
\[
  \text{gcd}(d^{n-1}x^n+y,\; d^{n-1}y^n+x)=e\le 2.
\]

Set
\[
  K=d^2xy+1.
\]
Then \(K\) is coprime to each of \(d\), \(x\), and \(y\). Choose \(n\) so that
\(n+1\) is a positive multiple of \(\varphi(K)\) and still large enough for the
previous identity to hold. By Euler's theorem,
\[
  d^{n+1}\equiv x^{n+1}\equiv y^{n+1}\equiv 1 \pmod K.
\]
Multiplying \(d^{n-1}x^n+y\) by \(d^2x\), we obtain
\[
  d^2x(d^{n-1}x^n+y)=d^{n+1}x^{n+1}+d^2xy \equiv 1+d^2xy \equiv 0 \pmod K.
\]
Since \(\gcd(d^2x,K)=1\), it follows that
\[
  K \mid d^{n-1}x^n+y.
\]
In the same way,
\[
  K \mid d^{n-1}y^n+x.
\]
Hence
\[
  K \mid \text{gcd}(d^{n-1}x^n+y,\; d^{n-1}y^n+x)=e.
\]
But \(e\le 2\), so \(K\le 2\). Since \(K=d^2xy+1\), this forces
\[
  d^2xy=1,
\]
and therefore
\[
  d=x=y=1.
\]
Thus \(a=dx=1\) and \(b=dy=1\).

So the unique solution is
\[
  (a,b)=(1,1).
\]
\end{proof}
